% =============================================================================
% Supplementary Section S4
% Multivariate analysis of the IllustrisTNG catalogue:
% environment as the primary independent variable
% Response to Reviewer 2, Major Comment 3 (revision v2, 2026-07)
% =============================================================================

\section*{S4. Multivariate analysis: environment as the primary variable}
\label{sec:S4}

\subsection*{Motivation}

The stratified analysis presented in the main text and in Section S3 defines the
Q- and R-dominated populations using \textit{internal} galaxy properties (gas
fraction and stellar mass) before examining their environmental dependence. By
construction, this design cannot by itself separate environmental effects
from population-dependent differences associated with morphology, quenching
history or baryonic subgrid physics. This section reports an analysis in which
the roles are reversed: environment is treated as the primary independent
variable and the internal properties are demoted to statistical controls.

\subsection*{Sample and method}

We use the TNG100-1 catalogue at $z = 0$, selecting galaxies with
$M_* > 10^9\,M_\odot$, giving $N = 53{,}363$ objects. Each galaxy contributes its
individual BTFR residual, computed exactly as for the observational samples
(vertical offset from a single global fit, no clipping, no outlier rejection).

Three predictors are used: the logarithmic local environmental density
$\rho \equiv \log \rho_{\mathrm{env}}$, the gas fraction
$f_{\mathrm{gas}} = M_{\mathrm{gas}}/(M_{\mathrm{gas}} + M_*)$, and $\log M_*$.
All three are standardised to zero mean and unit variance, so that fitted
coefficients are directly comparable in magnitude and the quadratic term is
evaluated about the sample mean. We fit three nested ordinary-least-squares
models with heteroskedasticity-consistent (HC3) standard errors:

\begin{align}
\text{M0:}\quad & \Delta \sim \rho + \rho^2 \\
\text{M1:}\quad & \Delta \sim \rho + \rho^2 + f_{\mathrm{gas}} + \log M_* \\
\text{M2:}\quad & \Delta \sim \rho + \rho^2 + f_{\mathrm{gas}} + \log M_*
  + \rho\!\cdot\!f_{\mathrm{gas}} + \rho\!\cdot\!\log M_* + \rho^2\!\cdot\!f_{\mathrm{gas}}
\end{align}

\subsection*{Results}

\begin{table}[htbp]
\centering
\caption{\textbf{Nested regression models, TNG100-1 ($N = 53{,}363$).} Standardised
predictors, HC3 robust standard errors. The quadratic environmental coefficient
remains positive and significant at every level, but is attenuated by roughly a
factor of eight between M0 and M2.}
\small
\begin{tabular}{@{}lrrr@{}}
\toprule
\textbf{Term} & \textbf{M0} & \textbf{M1} & \textbf{M2} \\
\midrule
Intercept                        & $-0.0790 \pm 0.0023$ & $-0.0285 \pm 0.0016$ & $-0.0683 \pm 0.0018$ \\
$\rho$                           & $-0.3092 \pm 0.0016$ & $-0.0226 \pm 0.0018$ & $+0.0323 \pm 0.0030$ \\
$\rho^2$                         & $+0.0790 \pm 0.0019$ & $+0.0285 \pm 0.0015$ & $+0.0095 \pm 0.0017$ \\
$f_{\mathrm{gas}}$               & ---                  & $+0.3509 \pm 0.0018$ & $+0.3680 \pm 0.0022$ \\
$\log M_*$                       & ---                  & $+0.0103 \pm 0.0010$ & $+0.0280 \pm 0.0011$ \\
$\rho \cdot f_{\mathrm{gas}}$    & ---                  & ---                  & $-0.1089 \pm 0.0033$ \\
$\rho \cdot \log M_*$            & ---                  & ---                  & $-0.0305 \pm 0.0013$ \\
$\rho^2 \cdot f_{\mathrm{gas}}$  & ---                  & ---                  & $+0.0467 \pm 0.0019$ \\
\midrule
$R^2$                            & 0.375   & 0.679  & 0.698 \\
AIC                              & 37313   & 1716   & $-1426$ \\
BIC                              & 37339   & 1761   & $-1355$ \\
\midrule
$\rho^2$ significance            & $42.6\sigma$ & $19.0\sigma$ & $5.7\sigma$ \\
\bottomrule
\end{tabular}
\label{tab:S4_nested}
\end{table}

\textbf{Survival of the quadratic term.} The environmental curvature remains
positive and formally significant in all three specifications
($p < 10^{-7}$ throughout). We draw attention, however, to the amplitude rather
than to the significance. Introducing gas fraction and stellar mass as controls
reduces the coefficient by approximately a factor of three, and adding the
interaction terms reduces it by a further factor of three; the coefficient in the
fully controlled model is therefore about eight times smaller than in the
environment-only fit. A substantial fraction of the raw environmental curvature
is thus attributable to the covariance between environment and internal galaxy
properties. Only a residual---robustly non-zero, but much reduced---component is
associated with environment once those properties are held fixed.

\textbf{Continuous sign inversion.} The interaction structure of M2 permits the
sign inversion to be examined without any prior subsample selection. The marginal
environmental slope at fixed stellar mass,
\begin{equation}
\frac{\partial \Delta}{\partial \rho}\bigg|_{\log M_*}
= \beta_\rho + \beta_{\rho f}\, f_{\mathrm{gas}},
\end{equation}
changes sign as a continuous function of gas content:

\begin{table}[htbp]
\centering
\caption{\textbf{Marginal environmental slope as a function of gas fraction (model M2).}}
\small
\begin{tabular}{@{}lrr@{}}
\toprule
\textbf{Population} & \textbf{$\partial \Delta / \partial \rho$} & \textbf{Significance} \\
\midrule
Gas-rich ($f_{\mathrm{gas}}$ at $+1\sigma$)      & $-0.0765 \pm 0.0024$ & $31.7\sigma$ \\
Star-dominated ($f_{\mathrm{gas}}$ at $-1\sigma$) & $+0.1412 \pm 0.0059$ & $24.1\sigma$ \\
\midrule
Interaction $\rho \cdot f_{\mathrm{gas}}$         & $-0.1089 \pm 0.0033$ & $33.0\sigma$ \\
\bottomrule
\end{tabular}
\label{tab:S4_marginal}
\end{table}

No population is defined in advance: the inversion emerges from the interaction
structure itself, as a continuous function of gas fraction.

\begin{figure}[htbp]
\centering
\includegraphics[width=\textwidth]{figureS4_multivariate.png}
\caption{\textbf{Multivariate analysis of the TNG100-1 catalogue with environment as
the primary variable.} \textbf{a}, The quadratic environmental coefficient across the
nested model hierarchy, with 95\% confidence intervals; the coefficient remains
positive throughout but is attenuated by roughly a factor of eight between M0 and
M2. \textbf{b}, The marginal environmental slope $\partial \Delta / \partial \rho$
from model M2, evaluated at fixed stellar mass, as a continuous function of gas
fraction, with its 95\% confidence band; the slope crosses zero, demonstrating the
sign inversion without any subsample pre-selection. \textbf{c}, Partial BTFR
residuals (gas fraction and stellar mass regressed out) binned by environmental
density, shown separately for the gas-rich upper quartile and the quenched lower
quartile; the two populations respond to environment in opposite directions.}
\label{fig:S4}
\end{figure}

\subsection*{Limitations}

We state the limitations of this analysis explicitly. The treatment controls for
gas fraction and stellar mass, but not for every galaxy property that co-varies
with environment in a hydrodynamical simulation---morphology, assembly and
quenching history, black-hole feedback state and satellite status among them. It
is a regression analysis and establishes association, not causality. The
attenuation of the curvature under controls should be read as evidence that a
meaningful fraction of the raw signal is population-driven; the surviving
component should be read as a lower bound on any genuinely environmental effect
within TNG, not as a measurement of a physical coupling.

We also note a feature of the catalogue relevant to the interpretation of the
quartile split in panel \textbf{c}: approximately 25\% of galaxies in the sample
have $f_{\mathrm{gas}} = 0$ exactly (fully quenched systems), so the lower
quartile consists of quenched galaxies rather than of a continuum of gas-poor
objects.

\subsection*{Reproducibility}

The analysis is implemented in \texttt{multivariate\_tng\_analysis.py}, archived
in the Zenodo deposit cited in the Code availability section of the main
manuscript. The script reads the derived TNG catalogue, fits the three nested
models, writes the full regression tables and generates
Figure~\ref{fig:S4} in a single execution.

